\expandafter\gdef\csname SNside1\endcsname{\node[anchor=north west,inner sep=0pt,text width=48mm,font=\SideFont,align=left] at ([xshift=138mm,yshift=-547.22pt]current page.north west) {\textcolor{SNink}{\textbf{最小反例 · 容量不等于表长}}\par\vspace{1.5mm}容量为100的数组可以表示空表，此时length=0；也可只保存一个元素。仅给数组上界不能推断有效元素数。插入判断必须同时检查位序是否合法与是否尚有空位。};}
\expandafter\gdef\csname SNside2\endcsname{\node[anchor=north west,inner sep=0pt,text width=48mm,font=\SideFont,align=left] at ([xshift=138mm,yshift=-283.73pt]current page.north west) {\textcolor{SNink}{\textbf{错项排除 · 随机存取不等于随机查值}}\par\vspace{1.5mm}顺序表给定位序后可用地址公式直接访问；只给目标值时，尚不知道它在哪个位置。无序表按值查找仍需逐个比较，不能把随机访问的O(1)套到所有查找上。};}
\expandafter\gdef\csname SNside3\endcsname{\node[anchor=north west,inner sep=0pt,text width=48mm,font=\SideFont,align=left] at ([xshift=138mm,yshift=-95.00pt]current page.north west) {\textcolor{SNink}{\textbf{错项排除 · 不搬移不等于常数时间}}\par\vspace{1.5mm}链表删除可只改链接，不需要移动其他元素的数据；但若还要从头定位前驱，操作总时间依然可能为线性。做题先拆分定位成本和局部结构修改成本。};}
\expandafter\gdef\csname SNside4\endcsname{\node[anchor=north west,inner sep=0pt,text width=48mm,font=\SideFont,align=left] at ([xshift=138mm,yshift=-95.00pt]current page.north west) {\textcolor{SNink}{\textbf{不变量 · 搬移保护未读区}}\par\vspace{1.5mm}插入向后搬移时，从后往前处理能确保源元素尚未被覆盖；删除向前搬移则从前往后读。记住需要保护的未读区，就能由目标下标推回正确循环方向。};}
\expandafter\gdef\csname SNside5\endcsname{\node[anchor=north west,inner sep=0pt,text width=48mm,font=\SideFont,align=left] at ([xshift=138mm,yshift=-95.00pt]current page.north west) {\textcolor{SNink}{\textbf{最小反例 · 两句交换会丢链}}\par\vspace{1.5mm}p原后继为r。若先写p-\textgreater next=s，再写s-\textgreater next=p-\textgreater next，第二句读到的是s，形成自环，r不再可从p到达。先暂存旧链接，是所有原地指针改写的基本纪律。\par\smallskip\begin{tikzpicture}[x=1mm,y=1mm,>=Stealth,node distance=5mm]
\node[draw=SNink!50,fill=SNblue,rounded corners=1mm,text width=34mm,align=center,inner sep=1.2mm,font=\fontsize{7.5}{9}\selectfont] (m0) at (0,0) {p};
\node[draw=SNink!50,fill=SNblue,rounded corners=1mm,text width=34mm,align=center,inner sep=1.2mm,font=\fontsize{7.5}{9}\selectfont] (m1) at (0,-11) {s};
\node[draw=SNink!50,fill=SNblue,rounded corners=1mm,text width=34mm,align=center,inner sep=1.2mm,font=\fontsize{7.5}{9}\selectfont] (m2) at (0,-22) {旧后继r};
\draw[->,SNink] (m0) -- (m1);
\draw[->,SNink] (m1) -- (m1);
\end{tikzpicture}
\par{\fontsize{7.5}{9}\selectfont 错误顺序使s自环，旧后继r不可达}};}
\expandafter\gdef\csname SNside6\endcsname{\node[anchor=north west,inner sep=0pt,text width=48mm,font=\SideFont,align=left] at ([xshift=138mm,yshift=-95.00pt]current page.north west) {\textcolor{SNink}{\textbf{不变量 · 固定间隔保持到尾}}\par\vspace{1.5mm}fast先走k步后，两指针始终相隔k条后继步。当fast到空位置时，slow距空位置也恰为k步，故指向倒数第k个结点。先走阶段遇空可直接识别k超过表长。};}
\expandafter\gdef\csname SNside7\endcsname{\node[anchor=north west,inner sep=0pt,text width=48mm,font=\SideFont,align=left] at ([xshift=138mm,yshift=-95.00pt]current page.north west) {\textcolor{SNink}{\textbf{最小反例 · 旧前驱必须还能取得}}\par\vspace{1.5mm}只有r、p两个相邻结点时，在p前插q的最终关系是r↔q↔p。若第一句覆盖p-\textgreater prior，之后再从它读取就失去r；四条目标链接全部成立才算插入完成。};}
\expandafter\gdef\csname SNside8\endcsname{\node[anchor=north west,inner sep=0pt,text width=48mm,font=\SideFont,align=left] at ([xshift=138mm,yshift=-95.00pt]current page.north west) {\textcolor{SNink}{\textbf{错项排除 · 字符相同不等于共享}}\par\vspace{1.5mm}两个结点都存字符i，也可能位于不同地址。公共后缀要求从某处起共享同一串结点，应比较指针而不是data；共享链不能再按两条独占链分别释放。};}
\expandafter\gdef\csname SNside9\endcsname{\node[anchor=north west,inner sep=0pt,text width=48mm,font=\SideFont,align=left] at ([xshift=138mm,yshift=-95.00pt]current page.north west) {\textcolor{SNink}{\textbf{不变量 · 头结点不随表空而消失}}\par\vspace{1.5mm}有效头结点始终存在，数据结点从L-\textgreater next开始。空表并不意味着L为空；统一前驱定位时pos=0对应头结点，才能让第一个位置的插入与其他位置共用同样两步链接。};}
\expandafter\gdef\csname SNside10\endcsname{\node[anchor=north west,inner sep=0pt,text width=48mm,font=\SideFont,align=left] at ([xshift=138mm,yshift=-95.00pt]current page.north west) {\textcolor{SNink}{\textbf{课程连接 · 表示影响访问局部性}}\par\vspace{1.5mm}连续数组通常能利用缓存中的相邻数据，分散结点可能增加缓存未命中。渐近复杂度相同也不保证实际用时相同；比较表示时同时检查数据规模、访问方式和存储开销。};}
\expandafter\gdef\csname SNside11\endcsname{\node[anchor=north west,inner sep=0pt,text width=48mm,font=\SideFont,align=left] at ([xshift=138mm,yshift=-95.00pt]current page.north west) {\textcolor{SNink}{\textbf{不变量 · 删后检查新后继}}\par\vspace{1.5mm}若当前前驱后连续两个结点都等于x，删掉第一个后马上前进，就可能跳过第二个。前驱只在保留结点时前进，才能让每个新后继都得到检查。};}
\expandafter\gdef\csname SNside12\endcsname{\node[anchor=north west,inner sep=0pt,text width=48mm,font=\SideFont,align=left] at ([xshift=138mm,yshift=-95.00pt]current page.north west) {\textcolor{SNink}{\textbf{最小反例 · 剩下候选仍需复核}}\par\vspace{1.5mm}序列1、2、3经抵消可留下3，但3只出现一次，未超过一半。抵消给出的是存在多数时的必要候选，并不证明多数存在；第二次计数补上充分性检验。};}
\AddToShipoutPictureFG{\begin{tikzpicture}[remember picture,overlay]\ifcsname SNside\thepage\endcsname\csname SNside\thepage\endcsname\fi\end{tikzpicture}}